
    %%%%%%%%%%%%%%%%%%%%%%%%%%%%%%%%%%%%%%%%%%%%%%%%%%%%%%%%%%%%%%%%%%%%%%%%%%%%%%%%
    %%% Classe do documento
    \documentclass[12pt, addpoints]{exam}
    \UseRawInputEncoding

    %%%%%%%%%%%%%%%%%%%%%%%%%%%%%%%%%%%%%%%%%%%%%%%%%%%%%%%%%%%%%%%%%%%%%%%%%%%%%%%%
    % preambulo.tex
    %
    % Arquivo de configuração de packages para uso o LaTeX, conforme minhas
    % preferências e modelos pessoais.
    %
    % Para maiores informações, visite:
    %    https://github.com/abrantesasf/latex
    %
    % NÃO ALTERE SE NÃO SOUBER O QUE ESTÁ FAZENDO!
    %%%%%%%%%%%%%%%%%%%%%%%%%%%%%%%%%%%%%%%%%%%%%%%%%%%%%%%%%%%%%%%%%%%%%%%%%%%%%%%%


    %%%%%%%%%%%%%%%%%%%%%%%%%%%%%%%%%%%%%%%%%%%%%%%%%%%%%%%%%%%%%%%%%%%%%%%%%%%%%%%%
    %%% Estruturas de controle:
    \input{static/utils/controles}

    %%%%%%%%%%%%%%%%%%%%%%%%%%%%%%%%%%%%%%%%%%%%%%%%%%%%%%%%%%%%%%%%%%%%%%%%%%%%%%%%
    %%% Compilação condicional em PDF:
    \input{static/utils/pdf}

    %%%%%%%%%%%%%%%%%%%%%%%%%%%%%%%%%%%%%%%%%%%%%%%%%%%%%%%%%%%%%%%%%%%%%%%%%%%%%%%%
    %%% Configurações de layout da página:
    \input{static/utils/layout}

    %%%%%%%%%%%%%%%%%%%%%%%%%%%%%%%%%%%%%%%%%%%%%%%%%%%%%%%%%%%%%%%%%%%%%%%%%%%%%%%%
    %%% Configurações para autores:
    \input{static/utils/autores}

    %%%%%%%%%%%%%%%%%%%%%%%%%%%%%%%%%%%%%%%%%%%%%%%%%%%%%%%%%%%%%%%%%%%%%%%%%%%%%%%%
    %%% Cabeçalho e rodapé:
    %%% Atenção! É MUITO DIFÍCIL ajustar apropriadamente os running headers
    %%% e running footers da primeira vez. Você pode confiar no LaTeX e deixar que
    %%% ele faça o ajuste automaticamente ou pode quebrar a cabeça e
    %%% tentar melhorar algo que o LaTeX já faz muito bem, mesmo que não fique
    %%% exatamente do jeito que você quer. O que você prefere? Se quiser ajustar
    %%% manualmente, descomente a linha abaixo e faça as configurações no arquivo
    %%% "utils/headings.tex".
    %\input{static/utils/headings}

    %%%%%%%%%%%%%%%%%%%%%%%%%%%%%%%%%%%%%%%%%%%%%%%%%%%%%%%%%%%%%%%%%%%%%%%%%%%%%%%%
    %%% Configuração para as fontes do documento:
    %%% Se você quiser usar fontes próprias ou do sistema, precisará ajustar as
    %%% configurações no arquivo "utils/fontes.tex".
    %%% ATENÇÃO: por padrão eu utilizo XeLaTeX com fontes proprietárias projetadas
    %%% por Matthew Butterick, adquiridas em https://mbtype.com, que não podem ser
    %%% distribuídas devido à licença de utilização. Se você usar XeLaTeX, você
    %%% DEVERÁ OBRIGATORIAMENTE ajustar as fontes no arquivo "utils/fontes.tex".
    \input{static/utils/fontes}

    %%%%%%%%%%%%%%%%%%%%%%%%%%%%%%%%%%%%%%%%%%%%%%%%%%%%%%%%%%%%%%%%%%%%%%%%%%%%%%%%
    %%% Sumário:
    \input{static/utils/toc}

    %%%%%%%%%%%%%%%%%%%%%%%%%%%%%%%%%%%%%%%%%%%%%%%%%%%%%%%%%%%%%%%%%%%%%%%%%%%%%%%%
    %%% Matemática:
    \input{static/utils/matematica}

    %%%%%%%%%%%%%%%%%%%%%%%%%%%%%%%%%%%%%%%%%%%%%%%%%%%%%%%%%%%%%%%%%%%%%%%%%%%%%%%%
    %%% Notas de rodapé:
    \input{static/utils/footnotes}

    %%%%%%%%%%%%%%%%%%%%%%%%%%%%%%%%%%%%%%%%%%%%%%%%%%%%%%%%%%%%%%%%%%%%%%%%%%%%%%%%
    %%% Referências cruzadas, links, citações:
    \input{static/utils/crossrefs}

    %%%%%%%%%%%%%%%%%%%%%%%%%%%%%%%%%%%%%%%%%%%%%%%%%%%%%%%%%%%%%%%%%%%%%%%%%%%%%%%%
    %%% Configurações para os hiperlinks em PDF:
    \input{static/utils/pdfconfig}

    %%%%%%%%%%%%%%%%%%%%%%%%%%%%%%%%%%%%%%%%%%%%%%%%%%%%%%%%%%%%%%%%%%%%%%%%%%%%%%%%
    %%% Glossário e índice remissivo:
    \input{static/utils/glossindex}

    %%%%%%%%%%%%%%%%%%%%%%%%%%%%%%%%%%%%%%%%%%%%%%%%%%%%%%%%%%%%%%%%%%%%%%%%%%%%%%%%
    %%% Referências bibliográficas:
    \input{static/utils/bibliografia}

    %%%%%%%%%%%%%%%%%%%%%%%%%%%%%%%%%%%%%%%%%%%%%%%%%%%%%%%%%%%%%%%%%%%%%%%%%%%%%%%%
    %%% Cores:
    \input{static/utils/cores}

    %%%%%%%%%%%%%%%%%%%%%%%%%%%%%%%%%%%%%%%%%%%%%%%%%%%%%%%%%%%%%%%%%%%%%%%%%%%%%%%%
    %%% Computação:
    \input{static/utils/computacao}

    %%%%%%%%%%%%%%%%%%%%%%%%%%%%%%%%%%%%%%%%%%%%%%%%%%%%%%%%%%%%%%%%%%%%%%%%%%%%%%%%
    %%% Gráficos:
    \input{static/utils/graficos}

    %%%%%%%%%%%%%%%%%%%%%%%%%%%%%%%%%%%%%%%%%%%%%%%%%%%%%%%%%%%%%%%%%%%%%%%%%%%%%%%%
    %%% Ícones:
    \input{static/utils/icones}

    %%%%%%%%%%%%%%%%%%%%%%%%%%%%%%%%%%%%%%%%%%%%%%%%%%%%%%%%%%%%%%%%%%%%%%%%%%%%%%%%
    %%% Blocos:
    \input{static/utils/blocos}

    %%%%%%%%%%%%%%%%%%%%%%%%%%%%%%%%%%%%%%%%%%%%%%%%%%%%%%%%%%%%%%%%%%%%%%%%%%%%%%%%
    %%% Boxes coloridos:
    \input{static/utils/colorbox}

    %%%%%%%%%%%%%%%%%%%%%%%%%%%%%%%%%%%%%%%%%%%%%%%%%%%%%%%%%%%%%%%%%%%%%%%%%%%%%%%%
    %%% Tabelas:
    \input{static/utils/tabelas}

    %%%%%%%%%%%%%%%%%%%%%%%%%%%%%%%%%%%%%%%%%%%%%%%%%%%%%%%%%%%%%%%%%%%%%%%%%%%%%%%%
    %%% Ambientes de listas:
    \input{static/utils/listas}

    %%%%%%%%%%%%%%%%%%%%%%%%%%%%%%%%%%%%%%%%%%%%%%%%%%%%%%%%%%%%%%%%%%%%%%%%%%%%%%%%
    %%% Ambientes floats e similares:
    \input{static/utils/floats}

    %%%%%%%%%%%%%%%%%%%%%%%%%%%%%%%%%%%%%%%%%%%%%%%%%%%%%%%%%%%%%%%%%%%%%%%%%%%%%%%%
    %%% Ajustes para a classe EXAM:
    \input{static/utils/exam}

    %%%%%%%%%%%%%%%%%%%%%%%%%%%%%%%%%%%%%%%%%%%%%%%%%%%%%%%%%%%%%%%%%%%%%%%%%%%%%%%%
    %%% Ajustes para textos:
    \input{static/utils/texto}

    %%%%%%%%%%%%%%%%%%%%%%%%%%%%%%%%%%%%%%%%%%%%%%%%%%%%%%%%%%%%%%%%%%%%%%%%%%%%%%%%
    %%% Ambientes, aliases, comandos e customizações em geral para serem
    %%% utilizadas nos documentos. Defina aqui qualquer customização específica
    %%% para seus documentos!
    \input{static/utils/customizacoes}

    %%%%%%%%%%%%%%%%%%%%%%%%%%%%%%%%%%%%%%%%%%%%%%%%%%%%%%%%%%%%%%%%%%%%%%%%%%%%%%%%
    %%% Ajuste do layout, espaçamento de linhas e etc.:
    \geometry{a4paper, portrait, left=2.0cm, right=2.0cm, top=2.0cm, bottom=2.0cm}
    \singlespacing

    %%%%%%%%%%%%%%%%%%%%%%%%%%%%%%%%%%%%%%%%%%%%%%%%%%%%%%%%%%%%%%%%%%%%%%%%%%%%%%%%
    %%% Configurações de cabeçalho e rodapé (não use fancyhdr pois ocorre conflito):
    \pagestyle{headandfoot}
    \runningheadrule
    \firstpageheader{}{}{}
    \runningheader{Sem disciplina}{}{Avaliação}
    \firstpagefooter{}{}{Página \thepage\ de \numpages}
    \runningfooter{}{}{Página \thepage\ de \numpages}
    \runningfootrule

    %%%%%%%%%%%%%%%%%%%%%%%%%%%%%%%%%%%%%%%%%%%%%%%%%%%%%%%%%%%%%%%%%%%%%%%%%%%%%%%%
    %%% Configurações para as propriedades do PDF:
    \hypersetup{
    hidelinks,           % Comente para web, descomente para impressão
    colorlinks=true,      % True para web, False para impressão
    pdftitle=AGo: Avaliação,
    pdfauthor=Matheus Endlich Silveira,
    pdfsubject=Sem disciplina,
    pdfkeywords=['Sem tag'],
    pdfinfo={
        CreationDate={}, % Ex.: D:AAAAMMDDHH24MISS
        ModDate={}       % Ex.: D:AAAAMMDDHH24MISS
    }
    }
    \input{static/utils/pdfconfig}

    %%%%%%%%%%%%%%%%%%%%%%%%%%%%%%%%%%%%%%%%%%%%%%%%%%%%%%%%%%%%%%%%%%%%%%%%%%%%%%%%
    %%% Começa o documento
    \begin{document}

    %%%%%%%%%%%%%%%%%%%%%%%%%%%%%%%%%%%%%%%%%%%%%%%%%%%%%%%%%%%%%%%%%%%%%%%%%%%%%%%%
    %%% Folha de instruções padronizadas da UVV para a prova
    \pdfbookmark[1]{Instruções}{instrucoes}
    \begin{figure}[H]
        \begin{center}
            \includegraphics[scale=0.3]{{static/imgs/logo-uvv.png}}
        \end{center}	
    \end{figure}

    \vspace{-1cm}

    \begin{table}[H]
        \centering
        \begin{tabular}{|llll|}
            \hline
            \multicolumn{2}{|l|}{\textbf{Disciplina:} Sem disciplina} & 
            \multicolumn{1}{l|}{\multirow{3}{*}{\textbf{\makecell{Nota:\\ \,\\ \,}}}} & 
            \multirow{3}{*}{\textbf{\makecell{Coordenador:\\ \,\\ \,}}} \\ 
            \cline{1-2}
            \multicolumn{2}{|l|}{\textbf{Professor:} Matheus Endlich Silveira \hspace{4.72cm} } & 
            \multicolumn{1}{l|}{} & \\ 
            \cline{1-2}
            \multicolumn{2}{|l|}{\textbf{Aluno:}} & 
            \multicolumn{1}{l|}{} & \\ 
            \hline
            \multicolumn{1}{|l|}{\textbf{Turma:} \hspace{6.3cm} } & 
            \multicolumn{1}{l|}{\textbf{Semestre:}} & 
            \multicolumn{2}{l|}{\textbf{Valor:} 10 pontos} \\ 
            \hline
            \multicolumn{1}{|l|}{\textbf{Data:}} & 
            \multicolumn{3}{l|}{\textbf{Avaliação:}} \\ 
            \hline
        \end{tabular}
    \end{table}

    \begin{center}
        \fbox{\setlength\fboxsep{0.3cm} \fbox{\parbox{15.4cm}{
            \begin{center}
                \textbf{INSTRUÇÕES PARA A REALIZAÇÃO DESTA AVALIAÇÃO:}
            \end{center}
            \begin{itemize}[leftmargin=*]
                \item Esta é a primeira avaliação da disciplina Sem disciplina, 
                    e engloba o conteúdo referente Sem tag.
                \item Esta avaliação contém 50 questões objetivas, \textbf{todas obrigatórias}.
                \item \textbf{Leia atentamente cada questão!} As questões objetivas terão uma,
                    e apenas uma única, resposta a ser marcada.
                \item \textbf{Desligue o celular} antes de começar. A avaliação é
                    \textbf{individual} e \textbf{sem consulta}.
                \item As questões podem ser respondidas, na prova, com lápis ou caneta. Ao final
                    da prova \textbf{{VOCÊ DEVERÁ PREENCHER O GABARITO COM CANETA
                    DE COR PRETA}} \textbf{\underline{OU AZUL}} para que seja feita a correção
                    automatizada através da leitura do padrão de respostas marcado no
                    gabarito.
                \item \textbf{CUIDADO ao preencher o gabarito} pois eles são \textbf{INDIVIDUAIS
                    e NOMEADOS} (cada estudante tem um gabarito próprio específico, com um
                    código de identificação único). \textbf{Questões rasuradas são anuladas}
                    pelo software de leitura do gabarito.
                \item Ao final da prova, \textbf{DEVOLVA A PROVA E O GABARITO} para o professor.
                \item Siga todas as normas de \textbf{Integridade Acadêmica} da disciplina.
                    Alunos que forem flagrados com qualquer espécie de ``cola'' ou trocando
                    informações com outros alunos terão suas avaliações recolhidas, as notas
                    zeradas, e a situação será encaminhada para a coordenação para a aplicação
                    das penalidades previstas pela universidade.
                \item Com 90 minutos de avaliação você terá tempo de até 2,min por questão,
                    em média.
                \item Boa avaliação!
            \end{itemize}
            \vspace{2cm}
        }}}
    \end{center}
    \newpage

    %%%%%%%%%%%%%%%%%%%%%%%%%%%%%%%%%%%%%%%%%%%%%%%%%%%%%%%%%%%%%%%%%%%%%%%%%%%%%%%%
    %%% Inicia o bloco de questões:
    \begin{questions}

    \newpage
    \question

Se \( \lim_{n \to \infty} \left( \frac{1}{n^2} + \frac{2}{n^2} + \cdots + \frac{n-1}{n^2} \right) = L \) então


\begin{choices}
\choice \( L = \infty \)
\choice \( L = 0 \)
\choice \( L = 2 \)
\choice \( L = 1/2 \)
\choice \( L = 1 \)
\end{choices}

\question

Seja a árvore binária abaixo a representação de um espaço de estados para um problema \( p \), em que o estado inicial é \( a \), e \( i \) e \( f \) são estados finais.



\begin{figure}[H]

    \centering

    \includegraphics[width=0.5\linewidth]{static/imgs/defaul.png}

    \caption{Questão 59}

    \label{fig:enter-label}

\end{figure}



Um algoritmo de busca em largura-primeiro forneceria a seguinte sequência de estados como primeira alternativa a um caminho-solução para o problema \( p \):
\begin{choices}
\choice \( a \ b \ d \ h \ e \ i \)
\choice \( a \ b \ d \ e \ f \)
\choice \( a \ b \ c \ d \ e \ f \)
\choice \( a \ b \ e \ i \)
\choice \( a \ c \ f \)
\end{choices}

\question

Assinale a proposição verdadeira


\begin{choices}
\choice Se \( x \) é um número real tal que \( x^2 \leq 4 \) então \( x \leq 2 \) e \( x \leq -2 \)
\choice Se \( x + y \) é um número racional então \( x \) e \( y \) são números racionais
\choice Se \( x \) e \( y \) são números reais tais que \( x < y \) então \( x^2 < y^2 \)
\choice Se \( x < -4 \) ou \( x > 1 \) então \( \frac{2x + 3}{x - 1} > 1 \)
\choice nenhuma das anteriores
\end{choices}

\question

Um sistema centralizado é um concentrador de recursos; um sistema distribuído apresenta seus recursos dispersos. Entretanto, nem todo conjunto de recursos computacionais dispersos pode ser considerado um sistema distribuído. Considerando um conjunto de computadores, assinale a alternativa que melhor corresponde às características necessárias para considerá-lo um sistema distribuído:
\begin{choices}
\choice Existência de memória compartilhada e relógios locais sincronizados 
\choice Existência de sistema operacional idêntico e hardware padronizado em todos os computadores
\choice Suporte de rede e um relógio global
\choice Suporte de rede e funções primitivas de comunicação 
\choice Existência de memória secundária compartilhada e protocolos de sincronização de estado
\end{choices}

\question

\begin{figure}[H]

    \centering

    \includegraphics[width=0.5\linewidth]{static/imgs/defaul.png}

    \caption{Enter Caption}

\end{figure}

Se \( O = (0, 0, 0) \); \( A = (2, 4, 1) \); \( B = (3, 1, 1) \) e \( C = (1, 3, 5) \) então o volume do sólido acima é


\begin{choices}
\choice 44
\choice 35
\choice 21
\choice 30
\choice 35/2
\end{choices}

\question

Um agente SNMP é um aplicativo que é executado:
\begin{choices}
\choice em um dispositivo de rede
\choice em computadores denominados de gerentes
\choice em roteadores com filtragem de pacotes com o objetivo de proteger acesso a rede
\choice a partir de um computador específico para monitorar a rede
\choice em "firewalls" com o objetivo de proteger acesso a rede
\end{choices}

\question

Sistemas de processamento de transações, tais como sistemas de reservas aéreas, devem prover um mecanismo que garanta que cada transação não é afetada por outras transações que possam estar ocorrendo ao mesmo tempo. Transações de duas fases obedecem a um protocolo que garante essa atomicidade. Em transações de duas fases:


\begin{choices}
\choice Verifica-se a disponibilidade de todas as travas antes de executar qualquer ação de travamento.
\choice Uma trava compartilhada sobre um objeto deve ser obtida antes de uma trava exclusiva sobre o objeto ser obtida.
\choice Qualquer objeto corretamente travado deve ser destravado antes que outro objeto possa ser travado.
\choice Todas as operações de leitura ocorrem antes da primeira operação de escrita.
\choice Todas as ações de travamento (\textit{lock}) ocorrem antes da primeira ação de destravamento.
\end{choices}

\question

O número de \textit{strings binárias} de comprimento 7 e contendo um par de zeros consecutivos é


\begin{choices}
\choice 94
\choice 92
\choice 91
\choice 90
\choice 95
\end{choices}

\question

Os protocolos de transporte atribuem a cada serviço um identificador único, o qual é empregado para encaminhar uma requisição de um aplicativo cliente ao processo servidor correto. Nos protocolos de transporte TCP e UDP, como esse identificador se denomina?
\begin{choices}
\choice Identificador do processo (PID)
\choice Porta
\choice Conexão
\choice Protocolo de aplicação
\choice Endereço IP
\end{choices}

\question

Três atletas \( A \), \( B \) e \( C \) competiram, ao pares, numa corrida de \( d \) metros. Considerando que cada atleta teve o mesmo desempenho (ou seja, a mesma velocidade) ao competir com adversários distintos, e sabendo-se que



\begin{itemize}

    \item \( A \) venceu \( B \) chegando 20 metros à frente

    \item \( B \) venceu \( C \) chegando 10 metros à frente

    \item \( A \) venceu \( C \) chegando 28 metros à frente,

\end{itemize}



podemos afirmar que a corrida tem
\begin{choices}
\choice 100 metros
\choice 50 metros
\choice 200 metros
\choice 110 metros
\choice 150 metros
\end{choices}

\question

(ENADE 2014) Considere o seguinte banco de dados ilustrado no diagrama relaconal

abaixo:

\begin{figure}[H]

  \centering

  \caption{Estados e fornecedores}

  \label{fig:estados}

  %\fbox{

    \includegraphics[width=0.5\linewidth]{static/imgs/defaul.png}

  %}\\

  %\footnotesize{Fonte: xxxx.}

\end{figure}

A expressão SQL que obtém os nomes dos estados para os quais não há fornecedores

cadastrados é:


\begin{choices}
\choice \begin{verbatim}SELECT e.nome

FROM estado e,

FROM fornecedor f

WHERE e.nome = f.uf;

\end{verbatim}


\choice \begin{verbatim}SELECT e.uf

FROM estado e

WHERE e.nome NOT IN (SELECT f.uf

                     FROM fornecedor f);

\end{verbatim}
\choice \begin{verbatim}SELECT e.nome

FROM estado e

WHERE e.uf IN (SELECT f.uf

               FROM fornecedor f);

\end{verbatim}


\choice \begin{verbatim}SELECT e.nome

FROM estado e

WHERE e.uf NOT IN (SELECT f.uf

                   FROM fornecedor f);

\end{verbatim}


\choice \begin{verbatim}SELECT e.nome

FROM estado e,

FROM fornecedor f

WHERE e.uf = f.uf;

\end{verbatim}


\end{choices}

\question

Uma definição lógica, abstrata e auto-contida dos objetos que modelam a

estrutura de armazenamento de dados, dos operadores que modelam o comportamento

e dos demais conceitos que modelam a integridade, e que, juntos, constituem a

máquina abstrata com a qual os usuários interagem é chamada de:


\begin{choices}
\choice Nenhuma das respostas acima
\choice Modelo relacional
\choice Modelo de dados
\choice Banco de dados
\choice Modelo entidade-relacionamento
\end{choices}

\question

Em relação aos metadados de um banco de dados, assinale a afirmação correta:
\begin{choices}
\choice São armazenados no catálogo (ou dicionário de dados) do sistema de gerenciamento de bancos de dados, que fica nas mesmas tabelas de dados dos usuários, permitindo assim que os metadados estejam sempre atualizados.
\choice São atualizados pelo SGBD, sempre que ocorre alguma operação de inserção, atualização ou remoção de dados.
\choice Como são sempre atualizados de forma automática pelo sistema de gerenciamento de bancos de dados, estão sempre prontos para nos dar uma ``foto' da estrutura do banco de dados em um momento no tempo.
\choice Armazenam informações sobre os dados dos usuários, ou seja, são dados de rotina dos usuários, aqueles dados que os usuários precisam para trabalhar no seu dia a dia.
\choice Armazenam informações como os dados de cadastro dos usuários, os valores de compra em uma ordem de compra, as turmas e disciplinas que os professores dão aula, e dados de recursos humanos.
\end{choices}

\question

Em uma lista circular duplamente encadeada com \( n \) elementos, o espaço ocupado apenas pelos apontadores é (assuma que um apontador ocupa \( p \) bytes):
\begin{choices}
\choice \( 4np \)
\choice \( np \)
\choice \( 2np \)
\choice \( 6np \)
\choice \( (np)^2 \)
\end{choices}

\question

Dentre as definições a seguir, ligadas ao conceito de normalização do modelo relacional, qual delas é \textbf{INCORRETA}?


\begin{choices}
\choice As formas normais se baseiam em certas estruturas de dependências.
\choice As relações que obedecem à primeira forma normal não apresentam anomalias.
\choice O objetivo da normalização é eliminar várias anomalias (ou aspectos indesejáveis) de uma relação.
\choice A normalização é um processo passo a passo reversível de substituição de uma dada coleção de relações por sucessivas coleções de relações as quais possuem uma estrutura progressivamente mais simples e mais regular.
\choice A primeira forma normal estabelece que os atributos da relação contêm apenas valores atômicos.
\end{choices}

\question

Dado um vetor \( u \in \mathbb{R}^2 \), \( u = (-3, 4) \), vamos denotar por \( v \) o vetor de \( \mathbb{R}^2 \) que tem tamanho 1 e é ortogonal a \( u \). Então \( v \) pode ser dado por


\begin{choices}
\choice \( \left( -\frac{4}{5}, \frac{2}{5} \right) \)
\choice \( \left( -\frac{4}{5}, -\frac{3}{5} \right) \)
\choice \( \left( -\frac{4}{5}, \frac{3}{5} \right) \)
\choice \( \left( -\frac{4}{5}, \frac{1}{5} \right) \)
\choice \( \left( \frac{3}{5}, \frac{4}{5} \right) \)
\end{choices}

\question

A derivada da função \( f(x) = x^x \) é igual a:
\begin{choices}
\choice \( x^x (\ln(x) + 1) \)
\choice \( x^x \)
\choice \( x^x (\ln(x) + x) \)
\choice \( x x^{x-1} \)
\choice \( x^x \ln(x) \)
\end{choices}

\question

Para cada \( n \in \mathbb{N} \), seja \( D_n = (0, 1/n) \), onde \( (0, 1/n) \) representa o intervalo aberto de extremos \( 0 \) e \( 1/n \). O conjunto diferença \( D_3 - D_{20} \) é igual a:


\begin{choices}
\choice \( (1/20, 1/3) \)
\choice \( D_{20} \cup D_3 \)
\choice \( D_{20} \)
\choice \( D_3 \)
\choice \( [1/20, 1/3) \)
\end{choices}

\question

São linguagens clássicas que representam aplicações científicas, aplicações 

comerciais, aplicações de inteligência artificial, e aplicações gerais,

respectivamente:
\begin{choices}
\choice R, COBOL, C, Python
\choice Fortran, COBOL, LISP, C
\choice Java, MATLAB, Scheme, JavaScript
\choice COBOL, Java, LISP, Fortran
\choice MATLAB, COBOL, SML, LISP
\end{choices}

\question

Considere \( C(x) \) uma função que define a complexidade de um problema \( x \); \( E(x) \) uma função que define o esforço (em termos de tempo) exigido para se resolver o problema \( x \). Sejam dois problemas denominados \( p1 \) e \( p2 \). Assinale a alternativa correta.


\begin{choices}
\choice \( C(p1 + p2) < C(p1) + C(p2) \)
\choice \( E(p1 + p2) < E(p1) + E(p2) \)
\choice Nenhuma das alternativas anteriores
\choice Se \( C(p1) < C(p2) \) então \( E(p1) < E(p2) \)
\choice Se \( C(p1) < C(p2) \) então \( E(p1) > E(p2) \)
\end{choices}

\question

Considere as seguintes tabelas em uma base de dados relacional:\\

\textbf{Departamento (CodDepto, NomeDepto)}\\

\textbf{Empregado (CodEmp, NomeEmp, CodDepto)}\\

Deseja-se obter uma tabela na qual cada linha é a concatenação de uma linha da tabela Departamento com uma linha da tabela de Empregado. Caso um departamento não possua empregados, sua linha no resultado deve conter vazio (NULL) nos campos referentes ao empregado. A operação de álgebra relacional que deve ser aplicada para combinar estas duas tabelas é:
\begin{choices}
\choice Projeção
\choice Junção interna
\choice Divisão
\choice União
\choice Junção externa
\end{choices}

\question

A sentença descritiva resultante de um processo de mensuração, correspondendo a

um fato bruto observado/conhecido que pode ser registrado, que não foi

processado para revelar seu significado/importância, que tem um certo

significado implícito e que deve estar corretamente formatado para o

armazenamento, processamento e apresentação é:
\begin{choices}
\choice Processamento
\choice Sabedoria
\choice Dado
\choice Informação
\choice Conhecimento
\end{choices}

\question

A medida da interconexão entre os módulos de uma estrutura de software é denominada e que também é usada em projetos orientados a objetos é: 
\begin{choices}
\choice acoplamento
\choice coesão
\choice ocultamento da informação
\choice unidade funcional
\choice abstração procedimental
\end{choices}

\question

Em um heap com \( n \) vértices existem:
\begin{choices}
\choice aproximadamente \( \log n \) folhas
\choice não mais que \( \left\lfloor \frac{n}{5} \right\rfloor \) folhas
\choice exatamente \( \left\lceil \frac{n}{5} \right\rceil \) folhas
\choice exatamente \( \left\lceil \frac{n}{2} \right\rceil \) folhas
\choice não menos que \( \frac{2n}{3} \) folhas
\end{choices}

\question

Na definição da aquisição de um novo software de

gerenciamento de bancos de dados (SGBD) para uma empresa da área de transporte 

coletivo urbano, a direção da área de informática conduziu o processo de decisão

da seguinte forma: foi designado um profissional da área de bancos de dados

(aquele com maior experiência na área) e atribuída a ele a tarefa de decidir

qual seria o melhor SBGD a ser adquirido. Esse profissional desenvolveu uma

série de estudos sobre as opções disponíveis, utilizando ténicas de simulação e

testes específicos, levando em conta que a quase totalidade dos dados que serão

armazenados são dados de cadastro altamente estruturados e que devem estar

disponíveis na Internet para consulta pelos usuários. Ao final, apresentou

ao diretor um relatório em que indicava claramente qual o melhor tipo de SGBD a

ser adquirido pela empresa. Com base nessa informação, o diretor da empresa

disparou o processo de compra do SGBD. Podemos concluir que o profissional de

bancos de dados escolheu um:
\begin{choices}
\choice SGBD Hierárquico, para armazenar corretamente e de forma otimizada todas os relacionamentos entre os ônibus, empregados e passageiros do transporte público.
\choice SGBD Grafo, para armazenar todos os relacionamentos entre os usuários do transporte público nos ônibus.
\choice SGBD Relacional, para armazenar de forma otimizada os dados que, em sua maioria, são altamente estruturados.
\choice Nenhuma das respostas acima.
\choice SGBD Não Relacional (NoSQL) já que, para armazenar todos os dados e ainda permitir a consulta pelos usuários na Internet, o melhor seria utilizar uma solução com bancos de dados não relacionais que são otimizados para uso na web.
\end{choices}

\question

Seu chefe que fazer uma campanha de venda para produtos que custam entre 20 e 30

reais, mas está preocupado porque ouviu de um dos vendedores que diversos

produtos nessa faixa de preço estão com o estoque zerado. Para verificar a

situação seu chefe pediu para que você prepare um relatório que mostre os

produtos que custam entre 20 e 30 reais e que estão com o estoque zerado. Você

preparou o seguinte relatório para seu chefe:

\begin{tcolorbox}

 \begin{verbatim}1  SELECT l.nome AS loja

2       , p.nome AS produto

3       , p.preco_unitario AS preco

4       , e.quantidade AS qtd

5  FROM produtos p

6  INNER JOIN estoques e ON (e.produto_id = p.produto_id)

7  INNER JOIN lojas l ON (e.loja_id = l.loja_id)

8  WHERE (p.preco_unitario >= 20 OR p.preco_unitario <= 30)

9    AND e.quantidade = 0

10 ORDER BY loja, produto;

 \end{verbatim}

\end{tcolorbox}

O seu relatório:
\begin{choices}
\choice Está errado pois a tabela que deveria estar na cláusula FROM é a tabela ``lojas', não a tabela ``produtos'.
\choice Nenhuma das alternativas anteriores.
\choice Está errado pois vai trazer também produtos fora da faixa de preço que o seu chefe quer.
\choice Está errado pois há um erro de sintaxe nas linhas 6 e 7.
\choice Está correto e trará exatamente, os produtos com o estoque zerado e na faixa de preço requisitada pelo seu chefe. Além disso ordena o resultado pelo nome da loja e pelo nome do produto.
\end{choices}

\question

Considere um problema em que são dados 5 objetos com os seguintes pesos e valores:\\

\text{pesos: } $(W_1, W_2, W_3, W_4, W_5) = (6, 10, 9, 5, 12)$\\

\text{valores: } $(P_1, P_2, P_3, P_4, P_5) = (8, 5, 10, 15, 7)$\\

Além disso, é dada uma mochila que suporta até 30 unidades de peso, para transportar os objetos. O objetivo do problema é preencher a mochila de tal forma que o valor total dos objetos a serem transportados seja o maior possível, mas sem exceder o limite de peso suportado pela mochila. Assuma que é permitido colocar fração de um objeto na mochila. Qual das seguintes alternativas corresponde ao valor máximo obtido no preenchimento da mochila:
\begin{choices}
\choice 43.1
\choice 21.5
\choice 12.2
\choice 30.34
\choice 38.83
\end{choices}

\question

Qual das seguintes opções \textbf{não é} uma vantagem do uso de um banco de

dados sobre o sistema de arquivos convencional?


\begin{choices}
\choice Estruturação de dados
\choice Eficiência
\choice Controle de acesso e visualização de dados
\choice Controle de redundância e inconsistência
\choice Independência de dados
\end{choices}

\question

Qual das seguintes afirmações sobre crescimento assintótico de funções não é verdadeira:
\begin{choices}
\choice Se \( f(n) = O(g(n)) \) então \( g(n) = O(f(n)) \)
\choice \( \log n^2 = O(\log n) \)
\choice \( 2n^2 + 3n + 1 = O(n^2) \)
\choice Se \( f(n) = O(g(n)) \) e \( g(n) = O(h(n)) \) então \( f(n) = O(h(n)) \)
\choice \( 2^{n+1} = O(2^n) \)
\end{choices}

\question

Considere as seguintes afirmações sobre redes neurais artificiais: \begin{enumerate}[label=\Roman*.] \item Um perceptron elementar só computa funções linearmente separáveis. \item Não aceitam valores numéricos como entrada. \item O "conhecimento" é representado principalmente através do peso das conexões. \end{enumerate} São corretas:
\begin{choices}
\choice I, II e III
\choice Apenas II e III
\choice Apenas I e II 
\choice Apenas I e III
\choice Apenas III
\end{choices}

\question

Sejam os seguintes predicados de uma linguagem de primeira ordem:



\begin{itemize}

    \item \( N(x) \): \( x \) é número;

    \item \( P(x) \): \( x \) tem propriedade \( P \);

    \item \( x < y \): \( x \) é menor que \( y \).

\end{itemize}



E sejam os símbolos:

\begin{itemize}

    \item \( \forall \): quantificador universal;

    \item \( \Rightarrow \): operador se-então;

    \item \( \neg \): operador de negação.

\end{itemize}



Para a fórmula: \( \forall x (N(x) \Rightarrow \neg \forall y (N(y) \Rightarrow y < x)) \), qual alternativa abaixo \textbf{NÃO} constitui uma tradução possível?


\begin{choices}
\choice Para todo número, existe um outro número que é maior do que ele.
\choice Não há um número menor do que outro número.
\choice Não há um número tal que todos os números são menores do que ele.
\choice Para todo número, não é verdade que qualquer número seja menor do que ele.
\choice Para qualquer \( x \), se \( x \) é número, então não é verdade que todos os números são menores do que ele.
\end{choices}

\question

Qual das afirmações a seguir, relativas à análise sintática, está INCORRETA?


\begin{choices}
\choice Uma das diferenças entre os diversos algoritmos de análise redutiva é a forma de identificar o \textit{handle} na pilha.
\choice Algoritmos de análise redutiva podem ser utilizados mesmo para gramáticas ambíguas.
\choice Analisadores sintáticos descendentes recursivos são mais simples de implementar do que analisadores sintáticos redutivos.
\choice As gramáticas LL podem descrever mais linguagens do que as gramáticas LR.
\choice Algoritmos descendentes recursivos podem ser utilizados para algumas gramáticas ambíguas.
\end{choices}

\question

Qual é o número mínimo de comparações necessário para encontrar o menor elemento de um conjunto qualquer não ordenado de \( n \) elementos?
\begin{choices}
\choice \( n \)
\choice 1
\choice \( n \log n \)
\choice \( n - 1 \)
\choice \( n + 1 \)
\end{choices}

\question

Dentre as definições a seguir, ligadas ao conceito de visões do modelo relacional, qual delas é \textbf{INCORRETA}?


\begin{choices}
\choice Uma visão relacional é uma relação virtual que nunca é materializada.
\choice O gerenciamento de visões envolve a conversão da consulta do usuário sobre as visões para a consulta sobre as relações base.
\choice Uma visão relacional é uma relação virtual, derivada de relações base a partir da especificação de operações da álgebra relacional.
\choice Programas aplicativos do banco de dados podem ser executados sobre visões de relações da base de dados.
\choice Uma visão é útil por representar uma percepção particular do banco de dados, compartilhado por muitos aplicativos.
\end{choices}

\question

Assinale a alternativa INCORRETA:


\begin{choices}
\choice Serviços orientados a conexão podem ser implementados em subredes que funcionam no modo datagrama.
\choice O controle de fluxo tem como objetivo garantir que nenhum dos parceiros de uma comunicação inunda o outro enviando pacotes mais rápido do que ele pode tratar.
\choice Os serviços orientados a conexão podem ajudar no controle de congestionamento através da diminuição da taxa de transmissão durante um congestionamento em andamento.
\choice Nos serviços orientados a conexões há a necessidade de estabelecimento de uma conexão antes da transferência dos dados.
\choice Os serviços orientados a conexões são sempre confiáveis garantindo a entrega ordenada e completa dos dados transmitidos.
\end{choices}

\question

Considere as seguintes afirmações sobre mecanismos de inferência em sistemas baseados em regras.

\begin{enumerate}[label=\Roman*.]

    \item O encadeamento regressivo tem pouca utilidade prática, pois deve partir do possível resultado.

    \item O encadeamento progressivo tanto pode ser em amplitude quanto em profundidade.

    \item Podem trabalhar com informações incertas ou incompletas. 

\end{enumerate}

São corretas:
\begin{choices}
\choice Apenas I e III
\choice Apenas III
\choice Apenas I e II
\choice I, II e III
\choice Apenas II e III
\end{choices}

\question

Considere as seguintes tabelas em uma base de dados relacional:\\

Departamento (\underline{CodDepto}, NomeDepto)\\

Empregado (\underline{CodEmp}, NomeEmp, CodDepto, Salario)\\

Considere a seguinte consulta SQL:

\begin{verbatim}

SELECT d.CodDepto, d.NomeDepto, SUM(e.Salario)

FROM Departamento d

JOIN Empregado e ON d.CodDepto = e.CodDepto

GROUP BY d.CodDepto, d.NomeDepto

HAVING COUNT(e.CodEmp) > 2 AND AVG(e.Salario) > 40;

\end{verbatim}



Qual o resultado da consulta?
\begin{choices}
\choice Para cada departamento que tem mais que dois empregados e cuja média salarial é maior que 40, obter um grupo de linhas que contém, para cada empregado do departamento, o código de seu departamento, seguido do nome de seu departamento, seguido da soma dos salários dos empregados do departamento.
\choice Para cada departamento que tem mais que dois empregados e cuja média salarial, considerando todos empregados do departamento, exceto os dois primeiros, é maior que 40, obter o código de departamento, seguido do nome do departamento, seguido da soma dos salários dos empregados do departamento.
\choice Para cada empregado que tem mais que dois departamentos, ambos com média salarial maior que 40, obter o código de departamento, seguido do nome do departamento, seguido da soma dos salários dos empregados do departamento.
\choice Para cada departamento que tem mais que dois empregados e cuja média salarial é maior que 40, obter o código de departamento, seguido do nome do departamento, seguido da soma dos salários dos empregados do departamento.
\choice A consulta não retorna nada pois está incorreta.
\end{choices}

\question

Em qual das situações abaixo um sistema de Raciocínio Baseado em Casos não deve ser utilizado?


\begin{choices}
\choice Quando especialistas conversam sobre seus domínios dando exemplos.
\choice Em aplicações de diagnóstico médico.
\choice Quando for fácil a obtenção de regras
\choice Quando a experiência for tão valiosa quanto o conhecimento em livros texto.
\choice Quando as regras utilizadas apresentam um grande número de exceções.
\end{choices}

\question

Sobre a técnica conhecida como Z-buffer é correto afirmar que:
\begin{choices}
\choice Nenhuma das alternativas acima está correta. 
\choice As primitivas geométricas precisam estar ordenadas de acordo com a distância em relação ao observador. 
\choice As dimensões do Z-buffer são independentes das dimensões do frame buffer. 
\choice É possível realizar o cômputo das variáveis envolvidas de forma incremental.
\choice É uma técnica muito comum de detecção de colisão. 
\end{choices}

\question

Qual das seguintes igualdades é verdadeira?



\begin{enumerate}

    \item[(I)] \( n^2 = \mathcal{O}(n^3) \)

    \item[(II)] \( 2 \cdot n + 1 = \mathcal{O}(n^2) \)

    \item[(III)] \( n^3 = \mathcal{O}(n^2) \)

    \item[(IV)] \( 3 \cdot n + 5 \cdot n \log n = \mathcal{O}(n) \)

    \item[(V)] \( \log n + \sqrt{n} = \mathcal{O}(n) \)

\end{enumerate}



As opções para a resposta correta são:
\begin{choices}
\choice somente I, II e V
\choice somente III, IV e V
\choice somente II, III e IV
\choice somente I, III e IV
\choice somente I e II
\end{choices}

\question Assinale a alternativa \textbf{incorreta}:
\begin{choices}
\choice O MariaDB é um SGBD
\choice O modelo relacional de dados foi criado por Edgar Frank Cood
\choice No modelo relacional, tudo o que o usuário percebe são tabelas
\choice Um SGBD não tem interface gráfica
\choice Modelo de dados é a mesma coisa que o diagrama relacional
\end{choices}

\question

Qual é a opção que descreve a tarefa executada pelo seguinte algoritmo escrito em Pascal? 

\begin{verbatim} 

procedure fazalgo (var x, var y)

begin 

	  x := x + y; 

	  y := x - y;

	  x := x - y; 

end 

\end{verbatim}
\begin{choices}
\choice divide y por x utilizando a subtração e retorna o resultado em x
\choice divide x por y utilizando a subtração e retorna o resultado em x
\choice calcula o mínimo múltiplo comum entre x e y e retorna o valor em x
\choice não altera os valores de x e y
\choice troca os valores de x e y
\end{choices}

\question

Considere as seguintes tabelas em uma base de dados relacional (chaves primárias sublinhadas):\\

Departamento (\underline{CodDepto}, NomeDepto)\\

Empregado (\underline{CodEmp}, NomeEmp, CodDepto)\\

Considere as seguintes restrições de integridade sobre esta base de dados relacional:\\

- Empregado.CodDepto é sempre diferente de NULL\\

- Empregado.CodDepto é chave estrangeira da tabela Departamento com cláusulas ON DELETE RESTRICT e ON UPDATE RESTRICT\\

Qual das seguintes validações não é especificada por estas restrições de integridade:
\begin{choices}
\choice Sempre que uma nova linha for inserida em Empregado, deve ser garantido que o valor de Empregado.CodDepto aparece na coluna Departamento.CodDepto.
\choice Sempre que uma nova linha for inserida em Departamento, deve ser garantido que o valor de Departamento.CodDepto aparece na coluna Empregado.CodDepto.
\choice Sempre que o valor de Departamento.CodDepto for alterado, deve ser garantido que não há uma linha com o antigo valor de Departamento.CodDepto na coluna Empregado.CodDepto.
\choice Sempre que o valor de Empregado.CodDepto for alterado, deve ser garantido que o novo valor de Empregado.CodDepto aparece em Departamento.CodDepto.
\choice Sempre que uma linha for excluída de Departamento, deve ser garantido que o valor de Departamento.CodDepto não aparece na coluna Empregado.CodDepto.
\end{choices}

\question

A velocidade de um ponto em movimento é dada pela equação\\

$v(t) = t e^{-0.01t} \, \text{m/s}$\\

O espaço percorrido desde o instante que o ponto começou a se mover até a sua parada total é


\begin{choices}
\choice \( 10^2 e^{-1} \, m \)
\choice \( 10^3 e^{-0.01} \, m \)
\choice \( 10^4 \, m \)
\choice \( 10^2 \, m \)
\choice \( (e^{-100} - 1) m \)
\end{choices}

\question

Redes semânticas, frames e lógica são formalismos utilizados principalmente em: 
\begin{choices}
\choice inferência em sistemas especialistas
\choice descoberta de conhecimento em bases de dados
\choice representação de conhecimento
\choice redes neurais
\choice IA distribuída
\end{choices}

\question

Que cláusula da SQL é fortemente relacionada com a lógica matemática?


\begin{choices}
\choice WHERE
\choice ORDER BY
\choice INNER JOIN
\choice SELECT
\choice FROM
\end{choices}

\question

A arquitetura ilustrada na figura abaixo é um exemplo típico de qual sistema?

(CU = unidade de controle; IS = \ingles{stream} de instruções; PU = unidade de

processamento; DS = \ingles{stream} de dados; LM = memória local)

\begin{figure}[H]

\begin{center}

\includegraphics[width=0.5\linewidth]{static/imgs/defaul.png}

\end{center}

\end{figure}

\vspace{-1cm}
\begin{choices}
\choice Processadores seqüenciais
\choice Multiprocessadores simétricos
\choice Acesso não uniforme à memória
\choice Processadores multicore
\choice Clusters
\end{choices}

\question

Avalie o banco de dados de professores e alunos, mostrado a seguir. Esse banco

de dados pretende armazenar todas as turmas, sendo que um professor pode dar

aula para vários alunos ao mesmo tempo, e um mesmo aluno pode ter aula com

diferentes professores ao mesmo tempo.

\begin{figure}[H]

  \centering

  \caption{Professores, turmas e alunos}

  \label{fig:professores}

  %\fbox{

    \includegraphics[width=0.5\linewidth]{static/imgs/defaul.png}

  %}\\

  %\footnotesize{Fonte: xxxx.}

\end{figure}

Esse banco de dados resolveráo problema de armazenar os dados das turmas?
\begin{choices}
\choice Não, pois o relacionamento entre 'turmas' e 'professores' está errado: a PK 'cod\_turma' deveria ter ido para a tabela 'professores' como uma FK, mas no diagrama está ao contrário.
\choice Não, pois as cardinalidades e os sentidos dos relacionamentos entre 'professores' e 'turmas', e entre 'turmas' e 'alunos' não cria um relacionamento N:N entre 'professores' e 'alunos'.
\choice Sim, pois as cardinalidades e os sentidos dos relacionamentos entre 'professores' e 'alunos' indica, no final, um relacionamento N:N entre essas duas tabelas.
\choice Não é possível determinar somente avaliando o diagrama exibido, seriam necessárias mais informações do dicionário de dados.
\choice Sim, pois o diagrama mostra claramente que um professor pode ter zero ou mais turmas, e cada aluno pode participar de zero ou mais turmas também.
\end{choices}

\question

Seja a seguinte linguagem, onde \(\epsilon\) representa o string vazio e \$ representa um marcador de fim de entrada:



$S \rightarrow ABCD \\$\\

$A \rightarrow a \mid \epsilon \\$\\

$B \rightarrow a \mid \epsilon \\$\\

$C \rightarrow c \mid \epsilon \\$\\

$D \rightarrow S \mid c \mid \epsilon$\\



É incorreto afirmar que:
\begin{choices}
\choice O conjunto FOLLOW(D) é igual a FOLLOW(S)
\choice O conjunto FIRST(A) = a, \(\epsilon\)
\choice O conjunto FOLLOW(B) = c, \$
\choice O conjunto FIRST(D) é igual ao conjunto FIRST(S)
\choice O conjunto FOLLOW(A) = a, c, \$
\end{choices}

\question

Quando dizemos que um banco de dados é em \ingles{cloud} e centralizado, estamos

nos referindo, respectivamente, à localização dos dados e da

infraestrutura. Verdadeiro ou falso?
\begin{choices}
\choice Não é possível determinar com os dados apresentados.
\choice Falso
\choice Verdadeiro
\end{choices}



    %%%%%%%%%%%%%%%%%%%%%%%%%%%%%%%%%%%%%%%%%%%%%%%%%%%%%%%%%%%%%%%%%%%%%%%%%%%%%%%%
    %%% Finaliza o bloco de questões:
    \end{questions}
    \end{document}
    